% Figura: campos con rotacional: F=-y ux + x uy (giro) y F=y ux (cizalla).
\begin{figure}[ht]
\centering
\begin{minipage}[t]{0.47\textwidth}\centering
\begin{tikzpicture}[x=0.58cm, y=0.58cm, baseline=(current bounding box.north)]
  \marcoD{-2.7}{2.7}{-2.7}{2.7}{-2,-1,0,1,2}{-2,-1,0,1,2}
  \ejesD{-2.7}{2.7}{-2.7}{2.7}
  \clipD{-2.7}{2.7}{-2.7}{2.7}{\flechas{-\y}{\x}{-2.4,-1.8,...,2.4}{-2.4,-1.8,...,2.4}{0.34}{3.4}}
\end{tikzpicture}\\[3pt]{\small $\vc{F}(\vc{r})=-y\,\ux+x\,\uy$}
\end{minipage}\hfill
\begin{minipage}[t]{0.47\textwidth}\centering
\begin{tikzpicture}[x=0.58cm, y=0.58cm, baseline=(current bounding box.north)]
  \marcoD{-2.7}{2.7}{-2.7}{2.7}{-2,-1,0,1,2}{-2,-1,0,1,2}
  \ejesD{-2.7}{2.7}{-2.7}{2.7}
  \clipD{-2.7}{2.7}{-2.7}{2.7}{\flechas{\y}{0}{-2.4,-1.8,...,2.4}{-2.4,-1.8,...,2.4}{0.34}{2.5}}
\end{tikzpicture}\\[3pt]{\small $\vc{F}(\vc{r})=y\,\ux$}
\end{minipage}
\caption{Campos con rotacional no nulo. (Izquierda) Giro rígido
$\vc{F}=-y\,\ux+x\,\uy$ (sentido antihorario): $\nab\times\vc{F}=2\,\uz$. Una
pequeña rueda colocada en cualquier punto gira con velocidad angular
constante. (Derecha) Cizalla $\vc{F}=y\,\ux$, de rotacional $-\uz$: no hay
líneas de corriente cerradas, pero la rueda gira igualmente (sentido horario)
porque el flujo es más rápido por un lado que por el otro.}
\label{fig:campos_rot}
\end{figure}
